% compile with xelatex or lualatex (the text contains Unicode math symbols)
\documentclass[11pt]{article}
\usepackage{amsmath,amssymb}
\usepackage{fontspec}
\setmainfont{DejaVu Serif}
\setmonofont{DejaVu Sans Mono}
\usepackage[margin=1in]{geometry}
\title{G (Catalan): continued fraction a(n) = 3n² + 9n + 7, b(n) = -2·n·(n + 1)\^{}3 equals (1) / (2 − 2·G (Catalan))}
\author{Euresis (automated discovery system)}
\date{\today}
\begin{document}
\maketitle
\paragraph{Status.} numerically verified, not proved.
\section{Method}
Found by a Ramanujan-Machine-style search over polynomial continued fractions; the fraction is evaluated with mpmath at two depths (to show convergence) and compared with the closed form. This is numerical evidence, not a proof.
\section{Evidence}
\begin{verbatim}
{
 "depths": [
  2000,
  4000
 ],
 "digits_agree": 130.3,
 "digits_converged": 130.0,
 "note": "numerical agreement is evidence, not proof"
}
\end{verbatim}
\section{Exact claim}
\begin{verbatim}
{
 "a": "3*n**2 + 9*n + 7",
 "b": "-2*n*(n + 1)**3",
 "digits": 100,
 "value": "(1) / (2 - 2*catalan)"
}
\end{verbatim}
\section{Novelty}
not in classical families or the Ramanujan Machine's published search domains; LIReC/OEIS/arXiv not checked.
\end{document}
